\documentclass[12pt, a4paper]{article}

% 引入必要的宏包
\usepackage[UTF8]{ctex}     % 中文支持
\usepackage{amsmath, amssymb, amsthm} % 数学符号和环境
\usepackage{geometry}       % 页面设置
\usepackage{enumitem}       % 列表定制
\usepackage{graphicx}       % 图片支持
\usepackage{fancyhdr}       % 页眉页脚

% 页面边距设置
\geometry{left=2.5cm, right=2.5cm, top=2.5cm, bottom=2.5cm}

% 宏定义：方便排版选择题选项
% 使用方法：\choices{选项A}{选项B}{选项C}{选项D}
\newcommand{\choices}[4]{
    \begin{enumerate}[label=\Alph*., itemsep=0pt, parsep=0pt, topsep=5pt, labelsep=0.5em]
        \item #1
        \item #2
        \item #3
        \item #4
    \end{enumerate}
}
% 横向排列的选项 (2列)
\newcommand{\choicesTwo}[4]{
    \begin{enumerate}[label=\Alph*., itemsep=0pt, parsep=0pt, topsep=5pt, labelsep=0.5em]
        \begin{minipage}{0.45\linewidth} \item #1 \end{minipage}
        \begin{minipage}{0.45\linewidth} \item #2 \end{minipage}
        \begin{minipage}{0.45\linewidth} \item #3 \end{minipage}
        \begin{minipage}{0.45\linewidth} \item #4 \end{minipage}
    \end{enumerate}
}
% 横向排列的选项 (4列)
\newcommand{\choicesFour}[4]{
    \begin{enumerate}[label=\Alph*., itemsep=0pt, parsep=0pt, topsep=5pt, labelsep=0.5em]
        \begin{minipage}{0.22\linewidth} \item #1 \end{minipage}
        \begin{minipage}{0.22\linewidth} \item #2 \end{minipage}
        \begin{minipage}{0.22\linewidth} \item #3 \end{minipage}
        \begin{minipage}{0.22\linewidth} \item #4 \end{minipage}
    \end{enumerate}
}

\title{\textbf{《高等数学》必刷习题——基础版}}
\author{}
\date{}

\begin{document}

\section*{第九章 多元微积分（基础）}

\begin{enumerate}
    \item 设二元函数 $ z = e^{xy} + xy $ ，则 $ \frac{\partial z}{\partial x}|_{(1,2)} = $ \underline{\hspace{1cm}}。
    \choicesFour{$ e^{2}+1 $}{$ 2e^{2}+2 $}{$ e+1 $}{$ 2e+1 $}

    \item 若可微函数 $ z = f(x, y) $ 在点 $ (x_{0}, y_{0}) $ 取得极小值，下列各项说法正确的是 \underline{\hspace{1cm}}。
    \choicesTwo
    {$ f(x_0,y) $ 在 $ y=y_{0} $ 处的导数大于 0}
    {$ f(x_0,y) $ 在 $ y=y_{0} $ 处的导数等于 0}
    {$ f(x_0,y) $ 在 $ y=y_{0} $ 处的导数小于 0}
    {$ f(x_0,y) $ 在 $ y=y_{0} $ 处的导数不存在}

    \item 对于函数 $ f(x,y) $ 的每一个驻点 $ (x_{0},y_{0}) $ ，令 $ A=f_{xx}(x_{0},y_{0}) $ ， $ B=f_{xy}(x_{0},y_{0}) $ ， $ C = f_{yy}(x_{0}, y_{0}) $ , 若 $ AC - B^{2} < 0 $ , 则函数 \underline{\hspace{1cm}}。
    \choicesFour{有极大值}{有极小值}{没有极值}{不定}

    \item 多元函数 $ f(x, y) $ 在点 $ (x_{0}, y_{0}) $ 处关于 $ y $ 的偏导数 $ f_{y}^{\prime}(x_{0}, y_{0}) = $ \underline{\hspace{1cm}}。
    \choicesTwo
    {$ \lim_{\Delta x \to 0} \frac{f(x_0 + \Delta x, y_0) - f(x_0, y_0)}{\Delta x} $}
    {$ \lim_{\Delta x \to 0} \frac{f(x_0 + \Delta x, y_0 + \Delta y) - f(x_0, y_0)}{\Delta x} $}
    {$ \lim_{\Delta y \to 0} \frac{f(x_0, y_0 + \Delta y) - f(x_0, y_0)}{\Delta y} $}
    {$ \lim_{\Delta x \to 0} \frac{f(x_0 + \Delta x, y_0 + \Delta y) - f(x_0, y_0)}{\Delta y} $}

    \item 若 $ f(x,y)=2x^{2}+y $ , 则 $ f_{x}^{'}= $ \underline{\hspace{1cm}}。
    \choicesFour{4}{0}{2}{-1}

    \item 求极限 $ \lim_{(x,y)\to(0,0)}\frac{\sqrt{xy+1}-1}{xy}= $ \underline{\hspace{2cm}}。

    \item 函数 $ z = \sqrt{y - x^{2} + 1} $ 定义域为 \underline{\hspace{2cm}}。

    \item 设 $ z = x^{y} $ ，则 $ \frac{\partial z}{\partial x} = $ \underline{\hspace{2cm}}。

    \item 设 $ f(x, y, z) = x^{y}y^{z} $ ，则 $ \frac{\partial f}{\partial y} = $ \underline{\hspace{2cm}}。

    \item 已知函数 $ z = x^{2} \arctan(2y) $ ，则全微分 $ dz = $ \underline{\hspace{2cm}}。

    \item 设 $ u = e^{\frac{x}{y}} $ , 求 $ \frac{\partial^{2}u}{\partial x \partial y} $。

    \item 求由方程 $ e^{z}-xyz=0 $ 所确定的二元函数 $ z=f(x,y) $ 的全微分。

    \item 求函数 $ f(x,y)=\frac{1}{3}x^{3}+2x-3xy+\frac{3}{2}y^{2} $ 的极值，并判断是极大值还是极小值。

    \item 设 $ z = z(x, y) $ 是由 $ x^{2}z + 2y^{2}z^{2} + y = 0 $ 确定的函数，求 $ \frac{\partial z}{\partial y} $。

    \item 已知 $ z = f(xy, 2x + 3y) $ ，其中 $ f(u, v) $ 具有连续偏导数，求 $ \frac{\partial z}{\partial x} $。
\end{enumerate}


\end{document}